\begin{abstract}
La detección temprana de anomalías en datos telemétricos es crítica para la seguridad y operación de sistemas espaciales. Este artículo presenta un framework basado en aprendizaje autosupervisado y modelos profundos para la detección automática de anomalías en telemetría de satélites y sistemas críticos. Se proponen arquitecturas híbridas que combinan Autoencoders Variacionales (VAE), redes contrastivas y reconstrucción temporal con mecanismos de atención. Las técnicas incluyen contrastive learning y augmentations específicas para series temporales multivariadas. Los modelos logran altas tasas de detección (F1-score >0.92) con baja tasa de falsos positivos en benchmarks como OPS-SAT y NASA SMAP/MSL. Se evalúa el rendimiento en entornos emulados edge, demostrando viabilidad para inferencia on-board. Se discuten trade-offs entre sensibilidad, generalización y robustez ante ruido y datos missing, junto con implicaciones para operaciones autónomas en misiones espaciales.
\end{abstract}